% ============================================================
%  附录B 介电函数的一般理论（原书印刷页 540–542）
%  底本：高等教育出版社《固体理论》第二版
%  （PDF 页 = 印刷页 + 20，原书 PDF 页 560–562）
%  公式编号自检：全附录 (B.1)–(B.15) 连续 15 个，无缺号，
%    与任务单 tags=15 一致。
%
%  已修正的 OCR/原书错误（除全角句点转"."、删 ~ 间隔符、
%  拆开行内公式与中文等常规项外）：
%  1. 章首"推迟格林函数 <<ρ_q|ρ_q^+>>_c ，"OCR 下标乱码，按
%     p540 复原为 <<ρ_q|ρ_q^+>>_ω。
%  2. (B.1) 整块 OCR 乱码按 p540 复原：1/ε(q,ω)=1−(4πe²/ħq²)
%     Σ_n|<n|ρ_q|0>|²(1/(ω+ω_{n0}+iη)−1/(ω−ω_{n0}+iη))；
%     其中 ω_{n0}=(1/ħ)(ℰ_n−ℰ_0)，花体 ℰ 统一 \mathscr。
%  3. 试探电荷密度式 OCR"−e£_q exp(…)"按 p540 复原为
%     −eξ_q exp(iq·r−iωt+ηt)（η=+0）。
%  4. (B.2) OCR"v(q)Σ_q ρ_q^+"按 p540 复原为 v(q)ξ_q ρ_q^+
%     （ξ_q 在哈密顿量中作常数因子，与 (B.3) 一致）。
%  5. (B.5) 与 (B.10)(B.7)(B.9)(B.12)(B.13)(B.15) 等 OCR 残缺
%     按 p540–542 复原；(B.5) 之首行 G_r(t)=−(i/ħ)θ(t)<[A(t),
%     B(0)]> 原书跨页印于 p540 页末，已并入本式编号块。
%  6. 原书 p541 (B.9) 为两行一块仅一号（∇·D 与 ∇·E），照排
%     （∇·E=4π×(试探电荷密度+感应电荷密度)）。
%  7. (B.9) 上下文"外扰动电荷密度"OCR 乱码按 p541 复原为
%     "试探电荷密度"。
%  8. 体例：< <A|B> > 统一 \langle\langle…\rangle\rangle；
%     虚单位 i、指数底 e 用正体；函数宗量 q 用粗体（照原书）。
%  遗留问题：
%  · 无。
% ============================================================
\chapter{介电函数的一般理论}

利用 $T=0$ K 时推迟格林函数 $\langle\langle\rho_{q}|\rho_{q}^{+}\rangle\rangle_{\omega}$ 的莱曼表示，可以证明电子体系介电函数的严格公式为
\begin{equation}
\begin{aligned}
\frac{1}{\varepsilon(\boldsymbol{q},\omega)}
&=1-\frac{4\pi e^{2}}{\hbar q^{2}}\sum_{n}
|\langle n|\rho_{q}|0\rangle|^{2}\times\\
&\quad\biggl(\frac{1}{\omega+\omega_{n0}+\mathrm{i}\eta}
-\frac{1}{\omega-\omega_{n0}+\mathrm{i}\eta}\biggr)
\end{aligned}
\end{equation}
其中 $\omega_{n0}=\dfrac{1}{\hbar}(\mathscr{E}_{n}-\mathscr{E}_{0})$，$|n\rangle$ 与 $\mathscr{E}_{n}$ 为 $H$ 的严格本征函数和能量本征值，$|0\rangle$ 为严格基态.

［证明］在互作用电子系统中人为的引入试探电荷作为外扰动，设试探电荷密度为
\begin{equation*}
-e\xi_{q}\exp(\mathrm{i}\boldsymbol{q}\cdot\boldsymbol{r}
-\mathrm{i}\omega t+\eta t)\qquad(\eta=+0)
\end{equation*}
$\xi_{q}$ 为一常数，则外扰动场的哈密顿量
\begin{equation*}
H_{\mathrm{e}}(t)=\sum_{j}\int\mathrm{d}\boldsymbol{r}\,
\frac{e^{2}\xi_{q}}{|\boldsymbol{r}_{j}-\boldsymbol{r}|}
\exp(\mathrm{i}\boldsymbol{q}\cdot\boldsymbol{r}
-\mathrm{i}\omega t+\eta t)
\end{equation*}
$\boldsymbol{r}_{j}$ 为电子坐标. 将 $e^{2}/|\boldsymbol{r}_{j}-\boldsymbol{r}|$ 作傅里叶展开，可得
\begin{equation}
\begin{aligned}
H_{\mathrm{e}}(t)&=\frac{4\pi e^{2}}{q^{2}}\xi_{q}\sum_{j}
\mathrm{e}^{\mathrm{i}\boldsymbol{q}\cdot\boldsymbol{r}_{j}}
\mathrm{e}^{-\mathrm{i}\omega t+\eta t}\\
&=v(q)\xi_{q}\rho_{q}^{+}\mathrm{e}^{-\mathrm{i}\omega t+\eta t}
\end{aligned}
\end{equation}
根据式(4.5.14)，外扰动的算符 $B$ 可取作
\begin{equation}
B=v(q)\xi_{q}\rho_{q}^{+}
\end{equation}
利用式(4.5.21)，外扰动所引起电子系统力学量 $A$ 的变化在线性响应范围为
\begin{equation*}
\langle A\rangle_{t}-\langle A\rangle_{t=-\infty}
=G_{\mathrm{r}}(\omega)\exp(-\mathrm{i}\omega t+\eta t)
\end{equation*}
若取 $A=\rho_{q}$，由于无外扰动时 $\rho_{q}$ 的统计平均值为零
\begin{equation*}
\langle\rho_{q}\rangle_{t=-\infty}=0
\end{equation*}
因此，系统密度起伏的 $q$ 分量为
\begin{equation}
\langle\rho_{q}\rangle_{t}=G_{\mathrm{r}}(\omega)
\exp(-\mathrm{i}\omega t+\eta t)
\end{equation}
其中 $G_{\mathrm{r}}(\omega)$ 是 $A=\rho_{q}$ 与 $B=\xi_{q}v(q)\rho_{q}^{+}$ 组成的推迟格林函数
\begin{equation}
\begin{aligned}
G_{\mathrm{r}}(t)&=-\frac{\mathrm{i}}{\hbar}\theta(t)
\langle[A(t),B(0)]\rangle\\
&=\xi_{q}v(q)\biggl\{-\frac{\mathrm{i}}{\hbar}\theta(t)
\langle[\rho_{q}(t),\rho_{q}^{+}]\rangle\biggr\}\\
&=\xi_{q}v(q)\langle\langle\rho_{q}(t);\rho_{q}^{+}\rangle\rangle
\end{aligned}
\end{equation}
其傅里叶变换为
\begin{equation}
G_{\mathrm{r}}(\omega)=\xi_{q}v(q)
\langle\langle\rho_{q}|\rho_{q}^{+}\rangle\rangle_{\omega+\mathrm{i}\eta}
\end{equation}
再利用莱曼表示式(A.14)
\begin{equation}
\begin{aligned}
\langle\langle\rho_{q}|\rho_{q}^{+}\rangle\rangle_{\omega+\mathrm{i}\eta}
&=\frac{1}{\hbar}\sum_{n}\biggl\{
\frac{\langle0|\rho_{q}|n\rangle\langle n|\rho_{q}^{+}|0\rangle}
{\omega-\omega_{n0}+\mathrm{i}\eta}-\\
&\qquad\frac{\langle0|\rho_{q}^{+}|n\rangle\langle n|\rho_{q}|0\rangle}
{\omega+\omega_{n0}+\mathrm{i}\eta}\biggr\}\\
&=\frac{1}{\hbar}\sum_{n}|\langle n|\rho_{q}|0\rangle|^{2}
\biggl(\frac{1}{\omega-\omega_{n0}+\mathrm{i}\eta}-\\
&\qquad\frac{1}{\omega+\omega_{n0}+\mathrm{i}\eta}\biggr)
\end{aligned}
\end{equation}
在后一等式中，利用了各向同性系统的对称关系
$|\langle n|\rho_{-q}|0\rangle|^{2}=|\langle n|\rho_{q}|0\rangle|^{2}$，于是
\begin{equation*}
\langle\rho_{q}\rangle_{t}=\xi_{q}v(q)
\langle\langle\rho_{q}|\rho_{q}^{+}\rangle\rangle_{\omega+\mathrm{i}\eta}
\exp(-\mathrm{i}\omega t+\eta t)
\end{equation*}
因此，外加试探电荷所引起的感应电荷密度为
\begin{equation}
-e\langle\rho_{q}\rangle_{t}\mathrm{e}^{\mathrm{i}\boldsymbol{q}
\cdot\boldsymbol{r}}=-e\xi_{q}v(q)
\langle\langle\rho_{q}|\rho_{q}^{+}\rangle\rangle_{\omega+\mathrm{i}\eta}
\exp(\mathrm{i}\boldsymbol{q}\cdot\boldsymbol{r}
-\mathrm{i}\omega t+\eta t)
\end{equation}
根据电动力学，电位移矢量 $\boldsymbol{D}$ 和电场 $\boldsymbol{E}$ 分别满足
\begin{align}
\nabla\cdot\boldsymbol{D}&=4\pi\times(\text{试探电荷密度})\nonumber\\
\nabla\cdot\boldsymbol{E}&=4\pi\times(\text{试探电荷密度}
+\text{感应电荷密度})
\end{align}
由于扰动和响应都按 $\exp(\mathrm{i}\boldsymbol{q}\cdot\boldsymbol{r}-\omega t+\eta t)$ 型式变化，因此 $\boldsymbol{D}(\boldsymbol{r},t)$ 与 $\boldsymbol{E}(\boldsymbol{r},t)$ 取下列形式
\begin{equation}
\begin{aligned}
\boldsymbol{D}(\boldsymbol{r},t)&=\boldsymbol{D}(\boldsymbol{q},\omega)
\exp(\mathrm{i}\boldsymbol{q}\cdot\boldsymbol{r}-\omega t+\eta t)\\
\boldsymbol{E}(\boldsymbol{r},t)&=\boldsymbol{E}(\boldsymbol{q},\omega)
\exp(\mathrm{i}\boldsymbol{q}\cdot\boldsymbol{r}-\omega t+\eta t)
\end{aligned}
\end{equation}
代入式(B.9)，求得 $\boldsymbol{D}$ 和 $\boldsymbol{E}$ 沿 $\boldsymbol{q}$ 方向的分量（纵向分量）满足方程
\begin{equation}
\mathrm{i}qD(q,\omega)=-4\pi e\xi_{q}
\end{equation}
\begin{equation}
\mathrm{i}qE(q,\omega)=-4\pi e\bigl[\xi_{q}+v(q)\xi_{q}
\langle\langle\rho_{q}|\rho_{q}^{+}\rangle\rangle_{\omega+\mathrm{i}\eta}
\bigr]
\end{equation}
纵向介电函数定义为
\begin{equation}
\varepsilon(q,\omega)=\frac{D(q,\omega)}{E(q,\omega)}
\end{equation}
由式(B.11)与式(B.12)立刻求得纵向介电函数与格林函数
$\langle\langle\rho_{q}|\rho_{q}^{+}\rangle\rangle_{\omega+\mathrm{i}\eta}$ 的关系为
\begin{equation}
\frac{1}{\varepsilon(q,\omega)}=1+v(q)
\langle\langle\rho_{q}|\rho_{q}^{+}\rangle\rangle_{\omega+\mathrm{i}\eta}
\end{equation}
将(B.7)代入上式，于是(B.1)得证.

根据(B.1)还得知，介电函数倒数的虚部为
\begin{equation}
\mathrm{Im}\biggl\{\frac{1}{\varepsilon(q,\omega)}\biggr\}
=\frac{\pi v(q)}{\hbar}\sum_{n}|\langle n|\rho_{q}|0\rangle|^{2}
\bigl[\delta(\omega+\omega_{n0})-\delta(\omega-\omega_{n0})\bigr]
\end{equation}
这个公式对于基态能的计算十分重要.
